% Version 0.28. Literal local Bloch--Kato and full CohomologyHyp installation.
% Pinned source: 7cbf7a0c32fca1badf05931879a37c3f8a2b03b4.
\section{Closing the local Bloch--Kato package on the original period models}
\label{sec:bloch-kato-installation}

\status{BX1--BX8 supply the full source predicate
\texttt{Published.\allowbreak Bloch\allowbreak Kato\allowbreak Local}, including both integral annihilator
directions and the original cocycle representatives.  Together with
LK, LD and DX this supplies the entire \texttt{Cohomology\allowbreak Hyp} record.
The classical inputs are Fontaine's standard period-ring construction
and the local theorem of Bloch--Kato for arbitrary de Rham
representations.  Their use is matched to the literal source models;
the global Bloch--Kato conjecture is not an input.}

\subsection{The interface that must actually be supplied}

For this chapter $\ell$ is any prime, including $2$, and
$K/\mathbb Q_\ell$ is an arbitrary finite extension inside the source's
specified algebraic closure.  Write $G_K$ for its absolute Galois group.
The input is an arbitrary \texttt{Loc\allowbreak Prime} $P$, an arbitrary
\texttt{Loc\allowbreak Ainf} $AI$ for it, and a witness
$T:\texttt{Loc\allowbreak Tied}(\ell,P,AI)$.  Thus the argument must work for every
such presentation of the period rings.  It need not, and does not,
claim a Bloch--Kato theorem for an abstract \texttt{Period\allowbreak Data} that has
no \texttt{Loc\allowbreak Tied} witness.

Put $Q=\mathbb Q_\ell$, $O=\mathbb Z_\ell$, and $PD=P.PD(K)$.
For a continuous finite-dimensional $Q$-representation $V$ the
source uses
\[
 \begin{split}
 Z^1(V)&=\{a:G_K\longrightarrow V:
     a\text{ continuous},\
     a(gh)=a(g)+g a(h)\},\\
 Z_f^1(PD,V)&=\{a\in Z^1(V):
     \exists x\in V\otimes_Q P.R.B,\
     a(g)\otimes1=gx-x\quad(g\in G_K)\},\\
 B^1(V)&=\{g\longmapsto gv-v:v\in V\}.
 \end{split}                                                   \tag{BX.1}
\]
Its \texttt{Finite\allowbreak Cohomology} is $Z_f^1/B^1$.  If $L\subset V$ is a
stable $O$-lattice, its finite lattice cocycles are
$Z_f^1(PD,V)\cap Z^1(G_K,L)$.  For $A=V/L$, its \texttt{Hf\allowbreak D} is the
sum of the image of $Z_f^1(PD,V)$ in functions to $A$ and the actual
coboundaries in $A$.  These are statements about representatives,
not merely unnamed cohomology subspaces.

Only two classical results specific to period rings are used below.
Fontaine's construction identifies the usual
$A_{\mathrm{cris}}$, $B_{\mathrm{cris}}$, $B_{\mathrm{dR}}^+$
and their filtration; see \emph{Le corps des p\'eriodes
$p$-adiques}, Ast\'erisque 223 (1994), \S1.5 and
\S2.3.3--2.3.4.
Bloch--Kato, \emph{$L$-functions and Tamagawa numbers of motives},
\S3, Proposition 3.8 and Corollary 3.8.4, give the finite-part
dimension formula and rational local annihilator theorem.
Both apply to every finite $K/Q$, every continuous de Rham $V$,
and every prime $\ell$.  Finite local Tate duality is the already
matched LK3/DX1 input.  The following proofs supply the intervening
identifications, lattice passage and record assembly.

% BEGIN BX1
\begin{proposition}[BX1: the actual tied period models and twist]
\label{claim:BX1}
The maps in \texttt{Loc\allowbreak Tied} identify $P.R.B$ and $P.R.BdR$,
as $Q$-algebras with $G_K$-action, with the standard
$B_{\mathrm{cris}}$ and $B_{\mathrm{dR}}$.
They carry the source filtration to
$\operatorname{Fil}^iB_{\mathrm{dR}}=t^iB_{\mathrm{dR}}^+$.
The character $PD.\chi$ is the actual $\ell$-adic cyclotomic
character on $G_K$.
\end{proposition}
\begin{proof}
The source's \texttt{Per} rings are explicit constructions.
For $C=\mathbb C_\ell$ they start with
$R=\varprojlim_{x\mapsto x^\ell}O_C/\ell$ and
$A_{\inf}=W(R)$, and the usual map
$\theta:A_{\inf}\to O_C$.
The ring called \texttt{Acris0} is the $A_{\inf}$-subalgebra of
$A_{\inf}[1/\ell]$ generated by
\[
                   \frac{x^n}{n!}\quad
          (x\in\ker\theta,\ n\ge0).
                                                               \tag{BX.2}
\]
Its $\ell$-adic completion is \texttt{Acris}.
This is the standard integral crystalline period ring: if
$\ker\theta=(\xi)$, then $x=a\xi$ gives
$x^n/n!=a^n\xi^n/n!$, so the displayed algebra is exactly
$A_{\inf}[\xi^n/n!:n\ge0]$.
This also identifies its specified action, induced on the tilt,
Witt vectors, divided powers and completion.

The source first inverts $\ell$ and then all logarithmic periods
of primitive compatible systems of $\ell$-power roots of unity.
This last localization is the usual inversion of a single $t$.
Indeed $T_\ell(\mu)$ is a free rank-one $O$-module: fixing a
primitive compatible system identifies every other such system
with its $u$-th power for a unique $u\in O^\times$.
The period logarithm is $O$-linear on this module, hence the
associated period is $ut$.  Its linearity follows from the
formal identity $\log(ab)=\log a+\log b$ on the divided-power
ideal and passage to compatible $\ell$-adic approximations.
All of these $u$ are units already before localization.
The two localizations therefore have inverse maps by their
universal properties.  This proves the claimed identification
of \texttt{Per.\allowbreak Bcris}.

Similarly, \texttt{Per.\allowbreak Bd\allowbreak Rp} is the completion of
$A_{\inf}[1/\ell]$ for powers of $\ker\theta$, and
\texttt{Per.\allowbreak Bd\allowbreak R} is its fraction field.  The completed kernel
is generated by $t$; its powers give $t^nB_{\mathrm{dR}}^+$.
The source's definition of negative filtration degrees uses the
quotient of the corresponding fractional ideals and thus gives
$t^iB_{\mathrm{dR}}^+$ for every integer $i$.
These are precisely the standard ring and filtration in the
classical statements just cited.

The supplied ring isomorphisms $T.eB$ and $T.eD$ commute with
the specified $O$-scalar maps by \texttt{z\allowbreak B}, \texttt{z\allowbreak D}.
They consequently commute with every $Q$-scalar: write
$q=\ell^{-m}z$, with $z\in O$, and use preservation of inverses.
The fields \texttt{act\allowbreak B}, \texttt{act\allowbreak D}, together with
$P.\texttt{PD\_act\allowbreak B}$ and $P.\texttt{PD\_actd\allowbreak R}$, give
$G_K$-equivariance.  The field \texttt{fil} gives the filtration
identification.  For the present endpoint these facts suffice;
no extra compatibility of a newly chosen ring embedding is
being inferred.

Finally let $g\in G_K$, and choose a primitive $\ell^m$-th root
$\zeta$.  The field $P.\texttt{chi\allowbreak Q\_cyc}$ gives an integer $a$
such that $g\zeta=\zeta^a$ and
$|\chi_Q(g)-a|_\ell\le\ell^{-m}$.
The first equality determines $a$ modulo $\ell^m$ as the
cyclotomic character.  The second identifies $\chi_Q(g)$ with
that residue modulo $\ell^m$.  For all $m$ these residues
determine a unique $\ell$-adic number, so
$\chi_Q(g)=\chi_\ell(g)$.  The equality $P.\texttt{PD\_chi}$
gives the asserted identity for $PD.\chi$.
No step here requires $\ell$ to be odd.
\end{proof}
% END BX1

% BEGIN BX2
\begin{proposition}[BX2: exact finite cocycles and period invariants]
\label{claim:BX2}
For every continuous finite-dimensional $Q$-representation $V$,
the source's $Z_f^1(PD,V)$ is the inverse image in $Z^1(V)$ of
the standard Bloch--Kato subspace $H_f^1(K,V)$.
Its \texttt{Dd\allowbreak R}, intersection with \texttt{Fil\allowbreak V(0)}, and
\texttt{Finite\allowbreak Cohomology} identify respectively with
$D_{\mathrm{dR}}(V)$, $\operatorname{Fil}^0D_{\mathrm{dR}}(V)$,
and $H_f^1(K,V)$ as $Q$-vector spaces.
In particular, \texttt{Is\allowbreak DR} is precisely the usual de Rham
condition.
\end{proposition}
\begin{proof}
Tensor the two $Q$-linear isomorphisms in BX1 with the identity
of $V$.  The resulting maps carry the source action
$g(v\otimes b)=gv\otimes gb$ to the standard diagonal action.
They therefore identify its common kernels of $g-1$ with the
usual period invariants.  The submodule \texttt{Fil\allowbreak V(0)} is the
image of $V\otimes_Q PD.\operatorname{Fil}^0$ in
$V\otimes_Q P.R.BdR$; its image is exactly
$V\otimes_Q B_{\mathrm{dR}}^+$.  Intersecting with invariants
gives the stated filtered subspace.

For each $a\in Z^1(V)$, the witness $x$ in (BX.1) is carried
bijectively to a standard $B_{\mathrm{cris}}$-coboundary witness.
Consequently the image of $[a]$ in
$H^1(K,V\otimes_Q B_{\mathrm{cris}})$ vanishes exactly when
$a\in Z_f^1(PD,V)$.  This is the definition of $H_f^1(K,V)$.
Every rational coboundary is finite, by the witness $v\otimes1$.
Taking the quotient by those actual coboundaries therefore
gives the asserted isomorphism on \texttt{Finite\allowbreak Cohomology}.
Continuity of $a$ is unchanged, since its coefficient space $V$
and topology have not changed.

The standard invariant space $D_{\mathrm{dR}}(V)$ is a
finite-dimensional $K$-vector space.  Its $Q$-dimension is
$[K:Q]\dim_KD_{\mathrm{dR}}(V)$.  Thus the source equality
\[
 \dim_Q\texttt{Dd\allowbreak R}(PD,V)=[K:Q]\dim_QV
                                                               \tag{BX.3}
\]
is exactly $\dim_KD_{\mathrm{dR}}(V)=\dim_QV$.
This verifies the precise antecedent of the classical theorem,
including the source's restriction of scalars.
\end{proof}
% END BX2

% BEGIN BX3
\begin{proposition}[BX3: the actual compact--discrete Tate passage]
\label{claim:BX3}
Let $L\subset V$ and $L'\subset V'=V^*(1)$ be dual stable
$O$-lattices, and put $A'=V'/L'$.
Then $H^1(K,L)$ is a finitely generated compact $O$-module,
$H^1(K,L)\otimes_OQ=H^1(K,V)$, and finite local Tate duality
induces a perfect continuous pairing
\[
 H^1(K,L)\times H^1(K,A')\longrightarrow Q/O.                    \tag{BX.4}
\]
The kernel of $H^1(K,V')\to H^1(K,A')$ is the image of
$H^1(K,L')$, a full lattice in $H^1(K,V')$.
The rational cup pairing
$H^1(K,V)\times H^1(K,V')\to Q$ is perfect and compatible
with (BX.4) by reduction modulo $O$.
\end{proposition}
\begin{proof}
Continuous cochains into $L$ are the inverse limit of continuous
cochains into $L/\ell^mL$.  At every finite quotient a cochain
has finite image, so it can be lifted to the next quotient by
choosing representatives on its finitely many clopen fibres.
The finite cohomology groups satisfy the Mittag--Leffler
condition: in each fixed finite group the descending images
stabilize.  The cocycle and coboundary exact sequences therefore
give
\[
 H^i(K,L)=\varprojlim_m H^i(K,L/\ell^mL).
                                                               \tag{BX.5}
\]
More explicitly, the possible obstruction is
$\varprojlim{}^1H^{i-1}(K,L/\ell^mL)$; the stabilized-image
property just proved makes it zero.  The usual proof chooses
successive cochain corrections after restricting to those
stable images.  In particular $M=H^1(K,L)$ is compact and
pro-$\ell$.

The long exact sequence for multiplication by $\ell$ embeds
$M/\ell M$ in the finite group $H^1(K,L/\ell L)$.
Choose finitely many lifts of its generators, and let $M_0$
be the closed $O$-submodule that they generate.  Then
$M/M_0=\ell(M/M_0)$.  A nonzero profinite pro-$\ell$ group has
a nonzero finite $\ell$-group quotient, on which multiplication
by $\ell$ is not surjective.  Hence $M/M_0=0$.
This proves finite generation and also identifies its compact
topology with its usual $O$-module topology.

A continuous cocycle on the compact group $G_K$ with values in
$V$ has bounded image and thus takes values in $\ell^{-m}L$
for some $m$.  A rational coboundary has the same property.
Localization of the cocycle quotient proves
$M\otimes_OQ=H^1(K,V)$.  The same argument applies to $L'$.
Continuous cochains into $A'$ have finite image and land in
some $\ell^{-m}L'/L'$, so
\[
 H^i(K,A')=\varinjlim_m H^i(K,\ell^{-m}L'/L').
                                                               \tag{BX.6}
\]
The two finite modules $L/\ell^mL$ and
$\ell^{-m}L'/L'$ are exact Tate duals with values in
$\ell^{-m}O/O(1)$.  LK3 and its normalized invariant identify
their cohomology pairings at each level.  The transition on the
left is reduction; its dual transition on the right is
inclusion.  Passing to (BX.5)--(BX.6) gives (BX.4), with its
original cup-product normalization.

The exact cohomology sequence of
$0\to L'\to V'\to A'\to0$ identifies the stated kernel.
Its image is the quotient of the finitely generated
$H^1(K,L')$ by its torsion and spans $H^1(K,V')$, so it is a
full lattice.  To check rational nondegeneracy without
localizing a torsion target, first repeat the finite-generation
argument in degree two.  Thus $H^2(K,L')$ is finitely generated
and the kernel of $H^2(K,L')\to H^2(K,V')$ is finite.
The same long exact sequence shows that
$N_0=\operatorname{im}(H^1(K,V')\to H^1(K,A'))$ has finite
index in $N=H^1(K,A')$.
By (BX.4), its annihilator in $M$ is finite.  Every nonzero
class in $M\otimes Q$ has, after scaling, a nontorsion
representative $m\in M$, so some element of $N_0$ pairs
nontrivially with $m$.  Lift that element to $H^1(K,V')$.
Compatibility of the cup formulas and the target map
$Q(1)\to Q/O(1)$ says that the nonzero value just obtained
is the reduction of the rational cup value.  Hence the
rational pairing has zero left kernel.  Interchanging
$L$ and $L'$ gives zero right kernel.  Both rational spaces
are finite-dimensional by the same finite-generation argument,
so the pairing is perfect.  This also proves the asserted
compatibility and normalization.
\end{proof}
% END BX3

% BEGIN BX4
\begin{proposition}[BX4: the literal finite-part dimension formula]
\label{claim:BX4}
For every representation in \texttt{Bloch\allowbreak Kato.\allowbreak Dim\allowbreak Formula}
whose source hypothesis \texttt{Is\allowbreak DR} holds,
\texttt{Finite\allowbreak Cohomology} is finite-dimensional and
\[
 \dim_Q H_f^1(PD,V)
   =\dim_Q V^{G_K}
      +\dim_Q D_{\mathrm{dR}}(PD,V)
      -\dim_Q\bigl(D_{\mathrm{dR}}(PD,V)
                    \cap\operatorname{FilV}(0)\bigr).          \tag{BX.7}
\]
Thus the complete source predicate
\texttt{Bloch\allowbreak Kato.\allowbreak Dim\allowbreak Formula}(PD) holds.
\end{proposition}
\begin{proof}
By BX2 the actual input representation is de Rham, with its
specified finite-dimensional module topology and continuous
action.  The finite-part formula of Bloch--Kato,
Corollary 3.8.4, applies:
\[
 \dim_QH_f^1(K,V)
      =\dim_QH^0(K,V)
          +\dim_Q\bigl(D_{\mathrm{dR}}(V)/
                        \operatorname{Fil}^0D_{\mathrm{dR}}(V)\bigr).
                                                               \tag{BX.8}
\]
It also gives finite dimensionality.  BX2 identifies the left
side and both period subspaces.  The source's \texttt{H0V} is
the intersection of the kernels of $\rho(g)-1$, hence exactly
$V^{G_K}$.  Since the filtered subspace is contained in the
finite-dimensional invariant space, the dimension of the
quotient in (BX.8) is the difference of the two dimensions
in (BX.7).  Interpreting that equality in $\mathbb Z$ gives
the source statement, whose subtraction is integer
subtraction.  This proof is uniform in $P,AI,T,K$ and in every
representation and topology quantified by \texttt{Dim\allowbreak Formula}.
\end{proof}
% END BX4

% BEGIN BX5
\begin{proposition}[BX5: both rational annihilator implications]
\label{claim:BX5}
The full predicate \texttt{Bloch\allowbreak Kato.\allowbreak Rational\allowbreak Duality}(PD)
holds for every $P,AI,T,K$ above and every source pairing satisfying
\[
                     \texttt{IsDual1}(PD,\rho,\rho',P_0).
\]
\end{proposition}
\begin{proof}
Keep the source's spaces $V,V'$ and its given bilinear form
$P_0:V\times V'\to Q$.  The map
\[
             V'\longrightarrow V^*(1),\qquad
             y\longmapsto\bigl(x\longmapsto P_0(x,y)\bigr)
                                                               \tag{BX.9}
\]
is injective by the second nondegeneracy condition and
surjective by the equality of dimensions in \texttt{Is\allowbreak Dual1}.
Its equivariance follows by substituting $g^{-1}x,y$ in
$P_0(gx,gy)=\chi_\ell(g)P_0(x,y)$, where BX1 has identified
the exact character.  It is a continuous linear isomorphism
for the specified finite-dimensional module topologies.
Thus this is an identification of the original pairing with
the evaluation pairing, not a replacement by an unrelated
duality.

For de Rham $V$, the rational local theorem of Bloch--Kato,
Proposition 3.8 with both local conditions equal to $f$,
says that $H_f^1(K,V)$ and $H_f^1(K,V^*(1))$ are exact
annihilators under the rational Tate cup pairing.
In particular it applies without a crystalline hypothesis.
The isomorphism (BX.9) preserves finite cocycles: tensor it
with the identity of $B_{\mathrm{cris}}$ and carry the
explicit coboundary witnesses in both directions.
BX2 therefore identifies the two subspaces with the source's
ones for $V$ and $V'$.

The source's cup cocycle is exactly
\[
              (g,h)\longmapsto P_0(a(g),g b(h)).
                                                               \tag{BX.10}
\]
Its \texttt{Cob\allowbreak L} predicate means that this cocycle is the
continuous coboundary of a $Q(1)$-valued $1$-cochain.
Consequently \texttt{Cob\allowbreak L} is equivalent to its class
vanishing in $H^2(K,Q(1))$.  The normalized Tate invariant is
an isomorphism from this group to $Q$, so vanishing is
equivalent to annihilation in the scalar pairing.
Every finite cohomology class has a finite cocycle
representative by BX2.  Testing all source finite cocycles
is therefore exactly testing all elements of the relevant
$H_f^1$.  Applying the two annihilator equalities gives
both implications and both orders in the literal definition
of \texttt{Rational\allowbreak Duality}.
\end{proof}
% END BX5

% BEGIN BX6
\begin{proposition}[BX6: integral annihilators, including torsion]
\label{claim:BX6}
For dual stable lattices $L,L'$ in de Rham $V,V'$,
define
\[
 \begin{split}
 H_f^1(K,L)&=\{m\in H^1(K,L):m_Q\in H_f^1(K,V)\},\\
 H_f^1(K,A')&=\operatorname{im}\bigl(
                H_f^1(K,V')\longrightarrow H^1(K,A')\bigr).
 \end{split}                                                   \tag{BX.11}
\]
These two subgroups are exact annihilators under (BX.4),
without a torsion-freeness assumption on $H^1(K,L)$.
\end{proposition}
\begin{proof}
Put $M=H^1(K,L)$, $U=M\otimes_OQ$, $W=H_f^1(K,V)$,
and $F=\{m\in M:m_Q\in W\}$.
The subgroup $F$ contains all torsion in $M$.  Moreover
$M/F$ is torsion-free: if $\ell^jm\in F$, then
$\ell^jm_Q\in W$, and $Q$-linearity of $W$ gives $m_Q\in W$.
By BX3 it follows that $M/F$ is a finite free $O$-module,
of rank $\dim_QU-\dim_QW$.
The subgroup $F$ is closed, since $W$ is closed in the
finite-dimensional $Q$-space $U$ and $M\to U$ is continuous.

The perfect pairing (BX.4) identifies its annihilator with
\[
       \operatorname{Hom}_{\mathrm{cont}}(M/F,Q/O)
            \simeq (Q/O)^{\dim U-\dim W}.                     \tag{BX.12}
\]
Let $W'=H_f^1(K,V')$, and let $B$ be its image in
$N=H^1(K,A')$.  BX5 gives $W'=W^\perp$ in the perfect
rational pairing.  Cup-product compatibility in BX3
therefore gives $B\subset F^\perp$.
The kernel of $W'\to N$ is the intersection of $W'$ with
the full lattice that is the image of $H^1(K,L')$.
This intersection is a full lattice in $W'$: multiply a
$Q$-basis of $W'$ by a sufficiently large common power of
$\ell$ to put it in the ambient lattice, and bound its
intersection above by a finite lattice in that basis.
Hence
\[
                     B\simeq (Q/O)^{\dim W'}.
                                                               \tag{BX.13}
\]
Rational duality gives
$\dim W'=\dim U-\dim W$.

For completeness, an inclusion of divisible $\ell$-primary
$O$-modules of the forms (BX.12)--(BX.13), with the same
finite corank, is equality.  A divisible module over the
PID $O$ is injective: extending a homomorphism from the
ideal $\ell^jO$ amounts to solving one equation
$\ell^jx=y$.  It follows that the inclusion splits.
Its complementary summand is divisible and $\ell$-primary
with zero corank.  Such a summand is zero, since a nonzero
element of order $\ell$ in a divisible module has
compatible preimages of arbitrarily high $\ell$-power
order, producing a copy of $Q/O$.  Thus $B=F^\perp$.

Finally the double annihilator of the closed subgroup $F$
under the perfect compact--discrete pairing is $F$.
Explicitly, if $m\notin F$, its nonzero image in the finite
free module $M/F$ survives in a finite quotient; a character
of that finite quotient into $Q/O$ detects it.
This gives $B^\perp=F$ and proves both directions, with the
torsion subgroup retained throughout.
\end{proof}
% END BX6

% BEGIN BX7
\begin{theorem}[BX7: the complete literal local Bloch--Kato predicate]
\label{claim:BX7}
For every prime $\ell$, every $P:\texttt{Loc\allowbreak Prime}(\ell)$,
every $AI:\texttt{Loc\allowbreak Ainf}(\ell,P)$ and every
$T:\texttt{Loc\allowbreak Tied}(\ell,P,AI)$, all three clauses
\[
 \texttt{Dim\allowbreak Formula}(PD),\qquad
 \texttt{Rational\allowbreak Duality}(PD),\qquad
 \texttt{Integral\allowbreak Duality}(PD)
                                                               \tag{BX.14}
\]
hold for every finite extension $K/Q$.
This supplies \texttt{Global\allowbreak Theory.\allowbreak Published.\allowbreak Bloch\allowbreak Kato\allowbreak Local}
with its original quantifiers and representatives.
\end{theorem}
\begin{proof}
BX4 and BX5 give the first two clauses.  To identify the
third, retain the given form $P_0$ and lattices $L,L'$ in
\texttt{Is\allowbreak Dual\allowbreak Lat1}.  The isomorphism (BX.9) sends $L'$
onto the exact $O$-dual of $L$: this is precisely the
source biconditional
\[
 y\in L'
   \quad\Longleftrightarrow\quad
    P_0(x,y)\in O\ \hbox{ for every }x\in L.                   \tag{BX.15}
\]
The target lattice $L_1$ in the scalar cyclotomic
representation has underlying submodule equal to the
image of $O\to Q$, by the explicit hypothesis in
\texttt{Integral\allowbreak Duality}.  Its quotient is therefore
the actual $Q/O(1)$, by BX1.

A lattice cocycle $d$ belongs to the source's
\texttt{lattice\allowbreak Finite\allowbreak Cocycles} exactly when its rational
class lies in $H_f^1(K,V)$, by BX2.
For a divisible cocycle $c:G_K\to V'/L'$, the definition
of \texttt{Hf\allowbreak D} says
\[
             c=q(a)+\delta\bar y
       \quad\text{for some }a\in Z_f^1(PD,V'),\
                                \bar y\in V'/L'.              \tag{BX.16}
\]
This is equivalent to $[c]$ lying in the image in (BX.11).
The forward implication is immediate.  For the reverse,
equality of the two classes says exactly that their
cocycle difference is an actual coboundary, yielding
(BX.16).  There is no need to lift $c$ to a rational
cocycle before testing this condition.

The source defines its integral cup using an arbitrary
choice of lift $\widetilde c(h)\in V'$:
\[
       (g,h)\longmapsto P_0(d(g),g\widetilde c(h))
                                  \pmod O.                    \tag{BX.17}
\]
Changing the lift changes this value by an element of $O$,
by stability of $L'$ and (BX.15).  Thus (BX.17) is the
ordinary cup product in $Q/O(1)$ and is independent of
the chosen section.  It is locally constant:
$c$ has finite image by compactness, hence is killed
by some $\ell^m$; the pairing then depends only on
$d$ modulo $\ell^mL$ and the action on the finite
module $\ell^{-m}L'/L'$.  Both are locally constant.
This argument also removes any continuity concern about
the arbitrary lift function itself.

The source's \texttt{Cob\allowbreak O} says exactly that (BX.17) is
in the group of locally constant $2$-coboundaries.
By the normalized invariant
$H^2(K,Q/O(1))\simeq Q/O$, this is exactly vanishing
of the Tate pairing in (BX.4).  BX6, applied to the
actual lattices and actual form, now says that $c$
satisfies (BX.16) if and only if it pairs trivially
with every finite lattice cocycle, and that a lattice
cocycle $d$ is finite if and only if it pairs trivially
with every such $c$.  These are exactly the two
conjuncts of \texttt{Integral\allowbreak Duality}.  All variables
were arbitrary in their original order, completing
the universally quantified predicate.
\end{proof}
% END BX7

% BEGIN BX8
\begin{theorem}[BX8: the full cohomology record]
\label{claim:BX8}
The complete predicates \texttt{Local\allowbreak Published\allowbreak Theorems}
and \texttt{Cohomology\allowbreak Hyp} are supplied by the reconstructed
local and global cohomology chains, with the classical
inputs stated in those chains.
\end{theorem}
\begin{proof}
The record \texttt{Local\allowbreak Published\allowbreak Theorems} has exactly
seven fields.  LK2, LK3 and LK4 give respectively
\texttt{lf\allowbreak K}, \texttt{lt\allowbreak K}, \texttt{le\allowbreak K}, on the original
finite local field and cyclotomic action.  DX1 gives
the native fields \texttt{lf}, \texttt{lt}, \texttt{le},
including nonfree finite modules and both degree-one
cup orders.  BX7 gives the entire field \texttt{bk}.
These seven statements have the same original
quantifiers, so they assemble the full local record.

The record \texttt{Cohomology\allowbreak Hyp} then has the following
complete field-by-field assembly:
\[
\begin{array}{c|l}
\text{field}&\text{supplier}\\ \hline
\texttt{loc}&\text{the seven-field record just assembled}\\
\texttt{pt}&\text{LD3, with the exact DX1 projection}\\
\texttt{gf}&\text{DX1: original global finite cohomology}\\
\texttt{ge}&\text{DX1: original truncated global Euler identity}\\
\texttt{uo}&\text{LD2: both unramified annihilator directions}\\
\texttt{gw}&\text{DX3: the complete Greenberg--Wiles formula}.
\end{array}                                                     \tag{BX.18}
\]
The normalized local invariant family is the one fixed
in LD3.  No choice of a different invariant is used to
combine the fields.  The global Euler field includes
the source's $2$-primary case through the truncated
degree-$0,1,2$ formula; its separate high-degree
vanishing clause and Greenberg--Wiles retain exactly
their source odd-prime antecedents.  The new local
Bloch--Kato field holds also at $\ell=2$.
Thus no field of the record is left as a new
bespoke input.
\end{proof}
% END BX8

\subsection{Consequence for the reconstruction frontier}

The older DX chapter correctly avoided assuming
\texttt{Bloch\allowbreak Kato\allowbreak Local} for its explicit local charts.
BX7 now supplies that larger predicate, and BX8
supplies the containing record for other consumers
that use it.  This is an additive completion of the
previous scope.  It does not by itself construct a
promodular support, an automorphic seed, or a uniform
automorphic catalogue.  Those source inputs have
their own arithmetic content and remain separately
accounted for in the dependency audit.
